Training an accurate object detection model requires large amounts of diverse images of the target object. Existing datasets for training are not always available, for example when the target is represented in a non-standard format or is uncommon. In those cases, the images must be captured manually, which is time-consuming and labour-intensive. Traditionally, one way to alleviate this problem has been to increase the number of training images by augmenting existing images.

In the last few years, text-to-image models capable of generating realistic synthetic images have developed rapidly. The SDXL text-to-image model from Stable Diffusion has become popular due to its open-source availability, enabling the creation of custom fine-tuned models that generate synthetic images. With only a small number of images, a generative model can be fine-tuned to produce synthetic images that match the content of the dataset. Ideally, these images are realistic enough to train an object detection model for use with real-world images. This work investigates whether such synthetic images can improve the performance of an object detection model.

In this work, an object detection model was developed to locate people in near-infrared (NIR) images in small indoor environments. To train the model, NIR-like synthetic images were generated using the SDXL text-to-image model, fine-tuned with 31 NIR images and the LoRA technique. These synthetic images were combined with real images to create a set of fine-tuned YOLOv8n object detection models, and they were ranged against each other.

The results showed that models fine-tuned solely with synthetic images outperformed the base model (AP\textsubscript{IoU:0.50}: 0.866-0.908 vs. 0.815). When real-world images were added to make up roughly one-tenth of the large synthetic dataset, the model's performance and stability increased even further. The best result was achieved by combining all 568 real-world images with 2500 synthetic images (AP\textsubscript{IoU:0.50}: 0.938). This work showed that synthetic images generated by a fine-tuned diffusion-based text-to-image model can improve the performance of object detection models to detect objects in real-world images.